% Paper paragraph for §8 Evaluation: Discriminative Power of the
% type-directed differential pipeline. Slot in between
% subsection{Bug-Finding Rate} (\S\ref{sec:eval-bugs}) and
% subsection{Coverage} (\S\ref{sec:eval-cov}). Numbers are from the
% Phase D run at $N = 1{,}000$ logged in
% \texttt{experiments/phase_d_mutants/outputs/summary.json}.

\subsection{Discriminative Power}
\label{sec:eval-discriminative}

The zero-disagreement headline of \S\ref{sec:eval-bugs} would be
epistemically empty if the pipeline could not detect a planted bug
either. We address this with a seeded-mutant study. Five
single-localised mutants were injected into \cedargo, each a
$\le\!4$-line edit reverting one correctness invariant: M1 swaps
every \texttt{Allow}/\texttt{Deny} return in
\texttt{Authorize}; M2 inverts the \texttt{\&\&} short-circuit branch
in \texttt{andEval}; M3 short-circuits \texttt{inEval} to
\texttt{true}; M4 swaps the \texttt{then} and \texttt{else} arms of
\texttt{ifThenElseEval}; M5 short-circuits \texttt{equalEval} to
\texttt{true}. Each mutant is a stand-alone unified diff; the
pipeline applies it, builds the harness, runs the §8 driver, and
reverts. All five compile without a workaround.
%
The 1000-tuple run is reported in
Table~\ref{tbl:phase-d-mutants}. M1 yields a $100.0\%$ disagreement
rate — every non-error decision is inverted, as predicted by
construction. The remaining four mutants yield disagreement rates
of $31.9\%$, $8.9\%$, $8.0\%$, and $6.0\%$ respectively, each
reflecting the share of the 25-shape \texttt{PolicyGen} corpus whose
final verdict depends on the mutated operator. The per-shape
breakdown isolates the responsible construct: M3 fires on
$89 / 89$ \texttt{permit$+$in} shapes (the only shapes that hit
the \texttt{in} evaluator on a load-bearing principal scope), M4 on
$80 / 80$ \texttt{permit$+$ite} shapes, and M5 on $60 / 60$
\texttt{permit$+$eq} shapes. No mutant yielded zero disagreements,
so the pipeline has no operator-blind spot among the mutated set.
M3, M4, and M5 do, however, illustrate how the disagreement rate is
upper-bounded by the share of policies whose decision actually
depends on the broken construct: an evaluator bug is invisible on a
policy whose final answer is fixed by other terms.

\begin{table}[t]
\centering
\footnotesize
\setlength{\tabcolsep}{4pt}
\begin{tabular}{llrr}
\toprule
ID & Injected bug                                & Disagr. & Rate \\
\midrule
M1 & flip \texttt{Allow}/\texttt{Deny}            & $1000 / 1000$ & $1.000$ \\
M2 & \texttt{\&\&} short-circuit on the wrong arm & $319 / 1000$  & $0.319$ \\
M3 & \texttt{in} unconditionally true             & $89 / 1000$   & $0.089$ \\
M4 & \texttt{ite} arms swapped                    & $80 / 1000$   & $0.080$ \\
M5 & \texttt{==} unconditionally true             & $60 / 1000$   & $0.060$ \\
\bottomrule
\end{tabular}
\caption{Phase D seeded-mutant study at $N = 1{,}000$ on the 25-shape
\texttt{PolicyGen} grammar. Each mutant is a single localised edit
to \cedargo (\texttt{$\le 4$ lines}) that reverts one correctness
invariant. The pipeline catches every mutant with a non-zero
disagreement rate; the rate scales with the share of generated
policies whose final decision depends on the mutated operator.
M1 lower-bounds the pipeline at $100\%$ when the mutation
is global; M3, M4, M5 illustrate the operator-coverage envelope of
the current 25-shape grammar.}
\label{tbl:phase-d-mutants}
\end{table}
